\documentclass{article}[12pt]
\usepackage{fullpage}
\usepackage{amsmath}

\begin{document}
\begin{center}
{\em  A Fourier-space four-field resonant response model for tokamak plasmas with applications to pedestal physics}\\[1ex]
by R.~Fitzpatrick and Y.~Lin\\[1ex]
{\bf Reply to Referees' Comments}\\[1ex]
~
\end{center}

First let us express our gratitude to the referees for their careful reading of our rather technical paper. Here is a detailed list of
our responses to the various comments. 

\begin{description}
\item[Referee 1:]

\begin{enumerate}
\item We have changed all references to Section~A to Appendix~A. 
\item We have added a comment to Sect.~IV.B that makes it clear that our list of devices that have successfully demonstrated
RMP-induced ELM suppression is not exhaustive. 
\item On p.~17, we have made it clear that the TM1 code has been benchmarked against both DIII-D and KSTAR experiments, and
have added some additional references. 
\item On p.~18, we have added a comment to the effect that the estimate $\chi_E\sim 1\,{\rm m^2/s}$ for ITER is actually very consistent
with the ITER89(y,2) scaling law. We have also made it clear that the estimate $\omega_{E\,0}=3\times 10^4\,{\rm rad./s}$ comes from
Fig.~8(b) of reference 55. This reference represents our best-guess estimate for the ITER rotation profile. As is clear from the figure, the
central plasma rotation is in the ion diamagnetic direction, and is of order a few tens of krad./s. 
\item We have replaced Ricatti by Riccati throughout the paper. We have inserted the missing ``of'' on p.~5. We have fixed the repetition of
``into" on p.~8. We have replaced ``imaginary curves" by ``dashed curves" in the caption to Fig.~6. We have inserted the missing ``is''
on p.~14. We have corrected the misspelling on RMP on p.~15. We have replaced a comma by a full-stop on p.~15. We have
replaced ``fourlayer" by ``four-field'' on p.~26. We have fixed the misspelling of ``of'' on p.~32. We have corrected the misspellings of
"Paz-Soldan" in the references. 
\end{enumerate}


\item[Referee 2:]

\begin{enumerate}
\item We have added language to Sect.~I.B to elucidate exactly what is meant by a four-field and a three-field layer response model. We have also
added language to Sect.~III.B to explain the physical significance of the ion sound parameter, $c_\beta$. 

\item We have given a reference to an explicit derivation of the cylindrical tearing mode equation, (1), in Sect.~II.B. We have also explained that it
is not consistent to include pressure gradient effects on tearing mode stability while ignoring the coupling of different poloidal harmonics. 

\item We have corrected the repetition of ``into'' on p.~8.

\item We have corrected the misspelling of ``RMP'' on p.~15.

\item On p.~20, we have clarified that the rational surface radius, $\hat{r}$, for $m=11$ corresponds to $q(\hat{r})=11$, whereas that for
$m=10$ corresponds to $q(\hat{r})=10$, and so on.

\item We have added comments to pp.~21, 22, and 23, pointing out the second decrease in ${\mit\Delta}_s$ (or ${\mit\Sigma}$) as $\hat{r}$ decreases, and
commenting that this behavior is the mechanism for island opening at the top of the pedestal. 

\item We have mentioned on p.~24 that the points at which the breakdown of shielding occurs are pushed radially inward from the flow-reversal
surface by an amount that seems to increase with increasing poloidal mode number. The origin of the poloidal mode number dependence is unclear. 

\end{enumerate}
\end{description}

\end{document}